\section{Kiểm thử hệ thống}

\subsection{Chiến lược và Phương pháp luận kiểm thử}

Nhằm đảm bảo hệ thống CoSpace vận hành với độ ổn định tuyệt đối, không có lỗi tiềm ẩn trong quá trình xử lý giao dịch tài chính và ngăn chặn hoàn toàn hiện tượng xung đột tài nguyên, nhóm áp dụng mô hình Kim tự tháp kiểm thử (Testing Pyramid) toàn diện từ thấp lên cao:

\begin{enumerate}
    \item \textbf{Kiểm thử đơn vị (Unit Testing):} Tập trung cô lập và thẩm tra tính đúng đắn của từng hàm, từng thuật toán nghiệp vụ độc lập (như tính điểm tương đồng Jaccard, thuật toán trừu tượng hóa khóa tương tranh, công thức khấu trừ phí hủy, mã hóa chữ ký số HMAC-SHA256). Lớp kiểm thử này sử dụng framework \textbf{JUnit 5 (Jupiter)} kết hợp thư viện giả lập đối tượng \textbf{Mockito 5}.
    \item \textbf{Kiểm thử tích hợp (Integration Testing):} Kiểm tra sự phối hợp mượt mà giữa các tầng phân lớp (Controller $\rightarrow$ Service $\rightarrow$ Repository), đảm bảo cơ chế bóc tách JWT của Spring Security và các câu lệnh truy vấn JPA giao tiếp chính xác với cơ sở dữ liệu.
    \item \textbf{Kiểm thử tích hợp cơ sở dữ liệu (Database Workflow Testing):} Thực thi trực tiếp các kịch bản kiểm thử SQL phức tạp để thẩm tra các ràng buộc toàn vẹn (\texttt{EXCLUDE USING gist}), hàm kích hoạt (Triggers), và các hàm khóa tương tranh \texttt{pg\_advisory\_xact\_lock}.
    \item \textbf{Kiểm thử đầu cuối (End-to-End System Testing):} Mô phỏng hành vi của người dùng thực tế trên giao diện, kiểm chứng toàn bộ luồng nghiệp vụ từ bước chọn chỗ ngồi trên sơ đồ SVG đến lúc quét mã QR Napas 24/7 và check-in tại quầy lễ tân.
    \item \textbf{Thẩm định yêu cầu phi chức năng (Non-Functional Verification):} Đánh giá độ trễ API, khả năng kiểm soát tương tranh và an toàn bảo mật.
\end{enumerate}

\subsection{Kiểm thử Đơn vị và Tích hợp trên tầng Back-end}

Toàn bộ hệ thống Back-end của CoSpace được bảo vệ bởi bộ kiểm thử tự động gồm \textbf{236 ca kiểm thử} (Unit \& Integration Tests) được viết bằng Java, đạt tỷ lệ vượt qua \textbf{100\% (236/236 PASS)}. 

\begin{table}[H]
\centering
\caption{Tổng hợp kết quả thực thi các bộ kiểm thử tự động Back-end}
\label{tab:backend-test-summary}
\renewcommand{\arraystretch}{1.25}
\begin{tabular}{|p{5.5cm}|p{2.2cm}|p{2cm}|p{2cm}|p{2.5cm}|}
\hline
\textbf{Bộ kiểm thử (Test Suite / Class)} & \textbf{Số ca test} & \textbf{Thành công} & \textbf{Thất bại} & \textbf{Tỷ lệ đạt} \\
\hline
\multicolumn{5}{|l|}{\textbf{1. Phân hệ Bảo mật, Định danh và Phân quyền (Security \& Auth)}} \\
\hline
\texttt{SecurityConfigRouteAuthorizationTest} & 21 & 21 & 0 & 100\% PASS \\
\texttt{SupabaseJwtAuthenticationConverterTest} & 11 & 11 & 0 & 100\% PASS \\
\texttt{BranchAccessGuardTest} & 18 & 18 & 0 & 100\% PASS \\
\texttt{JwtUtilTest} & 10 & 10 & 0 & 100\% PASS \\
\hline
\multicolumn{5}{|l|}{\textbf{2. Phân hệ Đặt chỗ và Vòng đời trạng thái (Booking \& Lifecycle)}} \\
\hline
\texttt{BookingServiceTest} & 28 & 28 & 0 & 100\% PASS \\
\texttt{BookingStateMachineTest} & 2 & 2 & 0 & 100\% PASS \\
\texttt{BookingExpirySchedulerTest} & 2 & 2 & 0 & 100\% PASS \\
\texttt{BookingLifecycleServiceTest} & 5 & 5 & 0 & 100\% PASS \\
\texttt{BookingAddonServiceTest} & 11 & 11 & 0 & 100\% PASS \\
\hline
\multicolumn{5}{|l|}{\textbf{3. Phân hệ Thanh toán và Đối soát VietQR (Payment \& PayOS)}} \\
\hline
\texttt{PaymentServiceTest} & 34 & 34 & 0 & 100\% PASS \\
\texttt{PayosServiceTest} & 17 & 17 & 0 & 100\% PASS \\
\texttt{MomoServiceTest} & 6 & 6 & 0 & 100\% PASS \\
\hline
\multicolumn{5}{|l|}{\textbf{4. Phân hệ Chính sách giá, Hủy phạt và Hoàn tiền (Pricing \& Refund)}} \\
\hline
\texttt{CancellationServiceTest} & 17 & 17 & 0 & 100\% PASS \\
\texttt{RefundServiceTest} & 7 & 7 & 0 & 100\% PASS \\
\texttt{PricingServiceTest} & 8 & 8 & 0 & 100\% PASS \\
\hline
\multicolumn{5}{|l|}{\textbf{5. Phân hệ Check-in, Lễ tân và Khách hàng thân thiết (Check-in \& Loyalty)}} \\
\hline
\texttt{CheckinServiceTest} & 13 & 13 & 0 & 100\% PASS \\
\texttt{LoyaltyRulesTest} & 11 & 11 & 0 & 100\% PASS \\
\texttt{AuthServiceTest} & 15 & 15 & 0 & 100\% PASS \\
\hline
\textbf{Tổng cộng toàn hệ thống Back-end} & \textbf{236} & \textbf{236} & \textbf{0} & \textbf{100\% PASS} \\
\hline
\end{tabular}
\end{table}

\subsubsection{Đặc tả ca kiểm thử Phân hệ Bảo mật và Phân quyền (Security \& Auth)}

\begin{table}[H]
\centering
\caption{Danh sách các ca kiểm thử tiêu biểu của Phân hệ Bảo mật và Phân quyền}
\label{tab:tc-security}
\renewcommand{\arraystretch}{1.25}
\begin{tabular}{|p{2.2cm}|p{3.5cm}|p{3.5cm}|p{4.2cm}|p{1.2cm}|}
\hline
\textbf{Mã ca test} & \textbf{Tên kiểm thử} & \textbf{Dữ liệu đầu vào (Input)} & \textbf{Kết quả mong đợi (Expected)} & \textbf{Kết quả} \\
\hline
TC-SEC-01 & Xác thực Bearer Token hợp lệ & Header chứa Supabase JWT hợp lệ & Nạp UserPrincipal vào SecurityContext, HTTP 200 & PASS \\
\hline
TC-SEC-02 & Từ chối Token đã hết hạn & JWT có \texttt{exp} trong quá khứ & Bị chặn bởi JwtFilter, trả về HTTP 401 Unauthorized & PASS \\
\hline
TC-SEC-03 & Chặn tài khoản bị vô hiệu hóa & JWT hợp lệ nhưng \texttt{status='suspended'} & Trả về HTTP 403 Forbidden kèm thông báo tài khoản bị khóa & PASS \\
\hline
TC-SEC-04 & Phân quyền Customer truy cập Admin API & User có Role \texttt{CUSTOMER} gọi \texttt{/api/admin/branches} & Bị chặn bởi \texttt{@PreAuthorize}, HTTP 403 Forbidden & PASS \\
\hline
TC-SEC-05 & Staff truy cập đúng chi nhánh & Staff chi nhánh A gọi cập nhật không gian chi nhánh A & Được phép thực thi (\texttt{hasBranchScope = true}), HTTP 200 & PASS \\
\hline
TC-SEC-06 & Chặn Staff thao tác chéo chi nhánh & Staff chi nhánh A gọi cập nhật không gian chi nhánh B & Ném ngoại lệ \texttt{AccessDeniedException}, HTTP 403 & PASS \\
\hline
TC-SEC-07 & Thẩm định cấu trúc JWT HS384 nội bộ & Token ký bằng Secret đối xứng 48-byte & Bóc tách chính xác userId, role, branchId & PASS \\
\hline
TC-SEC-08 & Kiểm tra chính sách CORS Whitelist & Origin từ domain lạ không nằm trong danh sách trắng & Preflight Request OPTIONS bị từ chối truy cập & PASS \\
\hline
\end{tabular}
\end{table}

\subsubsection{Đặc tả ca kiểm thử Phân hệ Đặt chỗ và Vòng đời trạng thái (Booking \& Lifecycle)}

\begin{table}[H]
\centering
\caption{Danh sách các ca kiểm thử tiêu biểu của Phân hệ Đặt chỗ}
\label{tab:tc-booking}
\renewcommand{\arraystretch}{1.25}
\begin{tabular}{|p{2.2cm}|p{3.5cm}|p{3.5cm}|p{4.2cm}|p{1.2cm}|}
\hline
\textbf{Mã ca test} & \textbf{Tên kiểm thử} & \textbf{Dữ liệu đầu vào (Input)} & \textbf{Kết quả mong đợi (Expected)} & \textbf{Kết quả} \\
\hline
TC-BOOK-01 & Tạo đặt chỗ hợp lệ đơn giản & Workspace khả dụng, thời lượng 4 giờ & Tạo đơn mới trạng thái \texttt{PENDING\_PAYMENT}, deadline 15p & PASS \\
\hline
TC-BOOK-02 & Áp dụng Rule \#33 giới hạn 3 đơn chờ & User đã có 3 đơn \texttt{PENDING\_PAYMENT} cố tạo đơn thứ 4 & Bị chặn bởi Advisory Lock, ném \texttt{IllegalStateException} & PASS \\
\hline
TC-BOOK-03 & Khóa tương tranh đồng thời trên Workspace & 2 luồng đồng thời xin đặt cùng 1 khung giờ & 1 luồng thành công (201), 1 luồng bị từ chối (409 Conflict) & PASS \\
\hline
TC-BOOK-04 & Tự động suy diễn Branch ID (Rule \#42) & Client gửi request không có branchId & Server tính toán chính xác từ \texttt{Workspace} $\rightarrow$ \texttt{Floor} $\rightarrow$ \texttt{Branch} & PASS \\
\hline
TC-BOOK-05 & Chặn đặt chỗ ngoài giờ mở cửa & Chi nhánh mở 08:00--22:00, đặt từ 23:00 & Bị từ chối với thông báo ngoài giờ vận hành & PASS \\
\hline
TC-BOOK-06 & Chặn đặt chỗ trùng lịch bảo trì & Workspace có lịch \texttt{MAINTENANCE} trùng khung giờ & Bị từ chối tạo đơn, bảo vệ tài nguyên đang sửa chữa & PASS \\
\hline
TC-BOOK-07 & Thêm dịch vụ tiện ích bổ sung & Booking hợp lệ + 2 ly cà phê espresso & Cập nhật snapshot giá dịch vụ và cộng dồn \texttt{addon\_amount} & PASS \\
\hline
TC-BOOK-08 & Hủy đơn quá hạn bởi Scheduler & Đơn \texttt{PENDING} vượt quá 15 phút & Chuyển sang \texttt{CANCELLED}, giải phóng ghế trên sơ đồ SVG & PASS \\
\hline
TC-BOOK-09 & Tự động chuyển đơn thành No-Show & Đơn \texttt{CONFIRMED} quá giờ \texttt{endAt} nhưng không check-in & Scheduler chuyển sang trạng thái \texttt{NO\_SHOW}, không hoàn tiền & PASS \\
\hline
TC-BOOK-10 & Gia hạn thời gian sử dụng (Extending) & Khách yêu cầu ngồi thêm 2 giờ kế tiếp & Kiểm tra không xung đột, cập nhật \texttt{endAt} và tạo Payment phụ & PASS \\
\hline
\end{tabular}
\end{table}

\subsubsection{Đặc tả ca kiểm thử Phân hệ Thanh toán VietQR và Webhook (Payment \& PayOS)}

\begin{table}[H]
\centering
\caption{Danh sách các ca kiểm thử tiêu biểu của Phân hệ Thanh toán VietQR}
\label{tab:tc-payment}
\renewcommand{\arraystretch}{1.25}
\begin{tabular}{|p{2.2cm}|p{3.5cm}|p{3.5cm}|p{4.2cm}|p{1.2cm}|}
\hline
\textbf{Mã ca test} & \textbf{Tên kiểm thử} & \textbf{Dữ liệu đầu vào (Input)} & \textbf{Kết quả mong đợi (Expected)} & \textbf{Kết quả} \\
\hline
TC-PAY-01 & Tạo liên kết thanh toán PayOS & OrderCode hợp lệ, số tiền 200,000 VNĐ & Trả về chuỗi mã VietQR động và đường dẫn thanh toán & PASS \\
\hline
TC-PAY-02 & Thẩm định chữ ký HMAC-SHA256 chuẩn & Webhook payload kèm chữ ký khớp với checksum key & Xác thực thành công (\texttt{verify = true}), tiếp tục xử lý & PASS \\
\hline
TC-PAY-03 & Chặn Webhook có chữ ký giả mạo & Webhook payload bị thay đổi số tiền hoặc signature sai & Từ chối xử lý, ghi log cảnh báo an ninh & PASS \\
\hline
TC-PAY-04 & Xử lý Webhook Idempotent (Chống lặp) & Gửi lại Webhook thành công lần thứ 2 cho cùng 1 giao dịch & Bỏ qua cập nhật trùng, trả về HTTP 200 OK ngay lập tức & PASS \\
\hline
TC-PAY-05 & Xử lý tiền vào trễ (Late Webhook Ghost) & Webhook báo thành công nhưng booking đã bị Scheduler hủy & Ghi nhận Payment paid, tự động tạo bản ghi hoàn tiền (Credit) & PASS \\
\hline
TC-PAY-06 & Xác nhận thanh toán tiền mặt tại quầy & Staff quầy nhận tiền mặt và bấm xác nhận & Cập nhật Payment \texttt{COMPLETED}, Booking \texttt{CONFIRMED} & PASS \\
\hline
TC-PAY-07 & Kích hoạt PayosProductionGuard & Chạy profile \texttt{prod} với cấu hình PayOS demo & Ngăn chặn khởi động máy chủ, bảo vệ an toàn tài chính & PASS \\
\hline
TC-PAY-08 & Bảo toàn bất biến tài chính Invariant & Đơn có chiết khấu và dịch vụ phát sinh & Đảm bảo \texttt{total = subtotal - discount + addon} tuyệt đối & PASS \\
\hline
\end{tabular}
\end{table}

\subsubsection{Đặc tả ca kiểm thử Phân hệ Chính sách giá, Hủy phạt và Hoàn tiền (Pricing \& Refund)}

\begin{table}[H]
\centering
\caption{Danh sách các ca kiểm thử tiêu biểu của Phân hệ Chính sách giá và Hoàn tiền}
\label{tab:tc-pricing}
\renewcommand{\arraystretch}{1.25}
\begin{tabular}{|p{2.2cm}|p{3.5cm}|p{3.5cm}|p{4.2cm}|p{1.2cm}|}
\hline
\textbf{Mã ca test} & \textbf{Tên kiểm thử} & \textbf{Dữ liệu đầu vào (Input)} & \textbf{Kết quả mong đợi (Expected)} & \textbf{Kết quả} \\
\hline
TC-PRI-01 & Tính giá thuê theo giờ phân tầng & Hot Desk, 4 giờ, bảng giá 30,000 đ/h & Tính ra subtotal = 120,000 VNĐ & PASS \\
\hline
TC-PRI-02 & Áp dụng chính sách giá riêng chi nhánh & Chi nhánh B có bảng giá ghi đè giá toàn cục & Ưu tiên lấy giá chi nhánh B thay vì giá mặc định & PASS \\
\hline
TC-PRI-03 & Hủy trước giờ bắt đầu $> 24$ giờ & Đơn đặt lúc 14:00 ngày mai, hủy trước 28h & Áp dụng chính sách hoàn 100\% tiền đã trả, phạt 0\% & PASS \\
\hline
TC-PRI-04 & Hủy trong khoảng $2 - 24$ giờ & Đơn hủy trước giờ bắt đầu 5 giờ & Áp dụng chính sách hoàn 50\% tiền đã trả, phạt 50\% & PASS \\
\hline
TC-PRI-05 & Hủy sát giờ $< 2$ giờ & Đơn hủy trước giờ bắt đầu 45 phút & Áp dụng chính sách hoàn 0\% (mất toàn bộ tiền cọc) & PASS \\
\hline
TC-PRI-06 & Không tìm thấy Policy hủy phù hợp & Khung giờ không khớp bất kỳ luật nào & Áp dụng Rule mặc định hoàn tiền 0\%, ghi log & PASS \\
\hline
TC-PRI-07 & Bảo lưu giá khi bảng giá cập nhật & Admin đổi giá phòng sau khi khách đã đặt & Giá trong đơn đặt chỗ cũ không bị thay đổi (Bảo toàn) & PASS \\
\hline
TC-PRI-08 & Staff hủy đơn do lỗi kỹ thuật chi nhánh & Phòng bị hỏng điều hòa đột xuất & Nhân viên chọn lý do chi nhánh, hoàn 100\% cho khách & PASS \\
\hline
\end{tabular}
\end{table}

\subsubsection{Đặc tả ca kiểm thử Phân hệ Lễ tân, Check-in và Khách hàng thân thiết (Check-in \& Loyalty)}

\begin{table}[H]
\centering
\caption{Danh sách các ca kiểm thử tiêu biểu của Phân hệ Lễ tân và Check-in}
\label{tab:tc-staff}
\renewcommand{\arraystretch}{1.25}
\begin{tabular}{|p{2.2cm}|p{3.5cm}|p{3.5cm}|p{4.2cm}|p{1.2cm}|}
\hline
\textbf{Mã ca test} & \textbf{Tên kiểm thử} & \textbf{Dữ liệu đầu vào (Input)} & \textbf{Kết quả mong đợi (Expected)} & \textbf{Kết quả} \\
\hline
TC-STF-01 & Check-in hợp lệ trong cửa sổ dung sai & Đến trước giờ bắt đầu 20 phút (trong $\pm 30$p) & Check-in thành công, trạng thái \texttt{CHECKED\_IN} & PASS \\
\hline
TC-STF-02 & Chặn Check-in quá sớm & Đến trước giờ bắt đầu 45 phút ($> 30$ phút) & Bị từ chối check-in kèm thông báo khung giờ cho phép & PASS \\
\hline
TC-STF-03 & Chặn Check-in quá muộn sau giờ kết thúc & Đến quầy sau khi ca làm việc đã kết thúc & Bị từ chối vì đơn đã bị đánh dấu \texttt{NO\_SHOW} & PASS \\
\hline
TC-STF-04 & Giới hạn duy nhất 1 check-in mở & Thử tạo 2 bản ghi check-in mở cho 1 booking & Ràng buộc \texttt{uq\_checkin\_open} chặn đứng lỗi trùng lặp & PASS \\
\hline
TC-STF-05 & Mở Running Tab gọi món tại chỗ & Khách đang ngồi gọi thêm 2 nước ngọt & Thêm dòng vào \texttt{booking\_services}, cập nhật tổng hóa đơn & PASS \\
\hline
TC-STF-06 & Check-out và thanh toán hóa đơn phát sinh & Khách trả phòng, còn nợ 50,000 đ tiền nước & Thu tiền mặt, đóng Running Tab, chuyển sang \texttt{COMPLETED} & PASS \\
\hline
TC-STF-07 & Khách thuê dài hạn Check-in hàng ngày & Hợp đồng thuê tháng (\texttt{is\_contract=true}) & Cho phép check-in/out nhiều lượt trên cùng 1 hợp đồng & PASS \\
\hline
TC-STF-08 & Tính toán thăng hạng thành viên Loyalty & User hoàn tất 10 đơn hoàn thành & Tự động nâng hạng Gold, áp dụng chiết khấu 10\% cho đơn kế tiếp & PASS \\
\hline
\end{tabular}
\end{table}

\subsubsection{Đặc tả ca kiểm thử Phân hệ Gợi ý Đối tác và Trợ lý ảo AI (Matching \& AI Assistant)}

\begin{table}[H]
\centering
\caption{Danh sách các ca kiểm thử tiêu biểu của Phân hệ Trí tuệ Nhân tạo}
\label{tab:tc-ai}
\renewcommand{\arraystretch}{1.25}
\begin{tabular}{|p{2.2cm}|p{3.5cm}|p{3.5cm}|p{4.2cm}|p{1.2cm}|}
\hline
\textbf{Mã ca test} & \textbf{Tên kiểm thử} & \textbf{Dữ liệu đầu vào (Input)} & \textbf{Kết quả mong đợi (Expected)} & \textbf{Kết quả} \\
\hline
TC-AI-01 & Tính toán chỉ số Jaccard kỹ năng chuẩn & 2 user có chung 3/5 kỹ năng & Điểm Jaccard đạt chính xác $3/5 = 0.60$ & PASS \\
\hline
TC-AI-02 & Cộng điểm cùng cơ sở chi nhánh (Rule \#18) & 2 user làm việc tại cùng cơ sở Quận 1 & Điểm tổng hợp được cộng thêm $+0.15$ điểm thưởng & PASS \\
\hline
TC-AI-03 & Chặn tự gợi ý kết nối chính mình & So khớp hồ sơ của chính người dùng hiện tại & Ràng buộc \texttt{check\_no\_self\_match} loại trừ bản thân & PASS \\
\hline
TC-AI-04 & Sinh câu lý do gặp gỡ bằng Gemini AI & Danh sách top ứng viên có điểm số $> 0.30$ & Mô hình sinh câu tóm tắt tương đồng tự nhiên bằng Tiếng Việt & PASS \\
\hline
TC-AI-05 & Cơ chế chuyển đổi mô hình dự phòng AI & Mô hình \texttt{gemini-3.5-flash} bị quá tải (429) & Tự động chuyển tiếp sang \texttt{gemini-3.7-flash} thành công & PASS \\
\hline
TC-AI-06 & Trả lời đàm thoại tư vấn chính sách phòng & Khách hỏi: "Tôi có được mang thú cưng vào không?" & Trợ lý ảo trích xuất đúng nội quy CoSpace và phản hồi thân thiện & PASS \\
\hline
\end{tabular}
\end{table}

\subsection{Kiểm thử Tích hợp Cấp Cơ sở dữ liệu (SQL Test Workflows)}

Bên cạnh các ca kiểm thử Java, nhóm xây dựng tập lệnh kiểm thử tự động mức cơ sở dữ liệu gồm \textbf{14 kịch bản chuyên sâu} trong tệp \texttt{database/tests/test\_full\_workflows.sql}. Các kịch bản này trực tiếp kích hoạt các tình huống biên phức tạp:

\begin{itemize}
    \item \textbf{Thử nghiệm vi phạm ràng buộc loại trừ GiST:} Cố tình chèn 2 bản ghi đặt chỗ có cùng \texttt{workspace\_id} và trùng lặp khoảng thời gian \texttt{tstzrange(start\_at, end\_at)}. Cơ sở dữ liệu lập tức chặn đứng câu lệnh thứ hai và ném ra lỗi mã hiệu PostgreSQL \texttt{23P01 (exclusion\_violation)}.
    \item \textbf{Thử nghiệm dung sai thời gian Check-in:} Kiểm chứng logic cho phép check-in sớm tối đa 30 phút và chặn các yêu cầu check-in quá trễ 30 phút sau khi ca làm việc kết thúc.
    \item \textbf{Thử nghiệm tính toán phí hoàn hủy theo 3 mức thời gian:} Kiểm tra hàm tính toán tỷ lệ hoàn tiền khi hủy trước 24 giờ (hoàn 100\%), hủy trong khoảng 2 -- 24 giờ (hoàn 50\%) và hủy dưới 2 giờ (hoàn 0\%).
    \item \textbf{Thử nghiệm bảo toàn tính độc lập giữa các chi nhánh (Branch Scoping Isolation):} Kiểm tra các tài khoản mang vai trò \texttt{BRANCH\_ADMIN} không thể truy vấn hoặc cập nhật dữ liệu của chi nhánh khác.
\end{itemize}

Kết quả thực nghiệm cho thấy 14/14 kịch bản SQL đều phản hồi kết quả chính xác theo đúng đặc tả hệ thống.

\subsection{Kịch bản kiểm thử tích hợp đầu cuối (End-to-End Scenarios)}

Để thẩm tra chất lượng tương tác giữa người dùng và hệ thống, nhóm xây dựng 5 kịch bản kiểm thử luồng nghiệp vụ hoàn chỉnh (E2E Test Cases) với đầy đủ các bước thực tế:

\subsubsection*{Kịch bản E2E 01: Đặt chỗ, Quét mã VietQR Napas 24/7 và Nhận phòng thành công}
\begin{itemize}
    \item \textbf{Mục tiêu:} Kiểm chứng luồng nghiệp vụ cơ bản chuẩn (Happy Path).
    \item \textbf{Các bước thực hiện:}
    \begin{enumerate}
        \item Người dùng đăng nhập vào hệ thống bằng tài khoản Google tại trang Đăng nhập (Hình~\ref{fig:ui_login_page}).
        \item Điều hướng tới trang Tìm kiếm không gian, chọn chi nhánh Quận 1, tầng 2 và nhấp chọn bàn làm việc \texttt{WD-01} trên sơ đồ tương tác SVG.
        \item Chọn khung giờ thuê từ 08:00 đến 12:00 ngày hôm sau và nhấn nút \textit{"Xác nhận đặt chỗ"}.
        \item Hệ thống hiển thị mã VietQR động cùng đồng hồ đếm ngược 15:00.
        \item Sử dụng ứng dụng Mobile Banking quét mã QR và chuyển khoản chính xác số tiền.
        \item Cổng PayOS gửi tín hiệu Webhook IPN về máy chủ CoSpace.
    \end{enumerate}
    \item \textbf{Kết quả mong đợi:} Trạng thái đơn đặt chỗ chuyển sang \texttt{CONFIRMED}, giao diện thanh toán tự động chuyển hướng sang màn hình Hóa đơn điện tử thông báo thành công.
    \item \textbf{Kết quả thực tế:} Hệ thống cập nhật trạng thái trong vòng 3.2 giây sau khi ngân hàng báo có. Đạt yêu cầu (\textbf{PASS}).
\end{itemize}

\subsubsection*{Kịch bản E2E 02: Xử lý tương tranh hai khách hàng cùng đặt một vị trí (Race Condition)}
\begin{itemize}
    \item \textbf{Mục tiêu:} Kiểm chứng cơ chế khóa \texttt{pg\_advisory\_xact\_lock} ngăn chặn hoàn toàn việc đụng lịch.
    \item \textbf{Các bước thực hiện:} Sử dụng công cụ tự động hóa gửi đồng thời hai yêu cầu đặt chỗ cho cùng một vị trí bàn làm việc \texttt{WD-02} tại cùng một khung thời gian trong vòng 50 mili-giây.
    \item \textbf{Kết quả mong đợi:} Duy nhất 1 yêu cầu gửi trước được cấp phép tạo đơn (\texttt{HTTP 201 Created}), yêu cầu gửi sau bị từ chối với thông báo lỗi rõ ràng (\texttt{HTTP 409 Conflict}).
    \item \textbf{Kết quả thực tế:} Cơ sở dữ liệu ghi nhận chính xác 1 bản ghi đặt chỗ, không xảy ra đụng lịch. Đạt yêu cầu (\textbf{PASS}).
\end{itemize}

\subsubsection*{Kịch bản E2E 03: Tự động thu hồi chỗ ngồi khi quá hạn thanh toán 15 phút}
\begin{itemize}
    \item \textbf{Mục tiêu:} Kiểm chứng hoạt động của tiến trình định kỳ \texttt{BookingExpiryScheduler}.
    \item \textbf{Các bước thực hiện:} Tạo một đơn đặt chỗ mới ở trạng thái \texttt{PENDING\_PAYMENT} nhưng cố tình không tiến hành thanh toán qua mã QR.
    \item \textbf{Kết quả mong đợi:} Sau khi hết 15 phút, tiến trình scheduler quét định kỳ 30s sẽ tự động đổi trạng thái đơn sang \texttt{CANCELLED}, đồng thời giải phóng vị trí bàn làm việc trên sơ đồ SVG trở về màu xanh khả dụng.
    \item \textbf{Kết quả thực tế:} Đúng 15 phút 18 giây, đơn được hủy tự động và vị trí bàn được mở lại cho các khách hàng khác. Đạt yêu cầu (\textbf{PASS}).
\end{itemize}

\subsubsection*{Kịch bản E2E 04: Quy trình Check-in tại quầy Lễ tân kết hợp Running Tab phát sinh}
\begin{itemize}
    \item \textbf{Mục tiêu:} Kiểm chứng chức năng vận hành thực tế của nhân viên quầy.
    \item \textbf{Các bước thực hiện:}
    \begin{enumerate}
        \item Khách hàng đến quầy và cung cấp mã đặt chỗ trước giờ bắt đầu 15 phút.
        \item Lễ tân nhập mã vào giao diện Check-in và nhấn \textit{"Check-in"}.
        \item Trong quá trình làm việc, khách hàng gọi thêm 2 ly cà phê và sử dụng dịch vụ in ấn. Lễ tân ghi nhận vào Running Tab của khách.
        \item Khi kết thúc buổi làm việc, lễ tân nhấn \textit{"Check-out \& In hóa đơn phát sinh"}.
    \end{enumerate}
    \item \textbf{Kết quả mong đợi:} Hệ thống cho phép check-in hợp lệ trong cửa sổ dung sai 30 phút, ghi nhận chính xác chi phí phát sinh và cập nhật trạng thái đơn sang \texttt{COMPLETED}.
    \item \textbf{Kết quả thực tế:} Hóa đơn tổng kết xuất chính xác từng danh mục dịch vụ đi kèm. Đạt yêu cầu (\textbf{PASS}).
\end{itemize}

\subsubsection*{Kịch bản E2E 05: Hủy đặt chỗ chủ động và Khấu trừ phí hoàn tiền theo chính sách}
\begin{itemize}
    \item \textbf{Mục tiêu:} Kiểm chứng tính đúng đắn của chính sách hoàn tiền tự động.
    \item \textbf{Các bước thực hiện:} Khách hàng có đơn đặt chỗ lúc 14:00 ngày mai. Vào lúc 10:00 sáng nay (còn cách giờ bắt đầu 28 tiếng), khách hàng nhấn nút \textit{"Hủy đặt chỗ"} trên trang Lịch sử (Hình~\ref{fig:ui_cancel_page}).
    \item \textbf{Kết quả mong đợi:} Hệ thống áp dụng quy tắc hủy trước 24 giờ, tính toán tỷ lệ hoàn tiền là 100\% giá trị đơn hàng, khởi tạo bản ghi hủy và giải phóng lịch trình.
    \item \textbf{Kết quả thực tế:} Giao diện hiển thị số tiền hoàn trả 100\%, lịch sử đặt chỗ cập nhật trạng thái \texttt{CANCELLED}. Đạt yêu cầu (\textbf{PASS}).
\end{itemize}

\subsection{Đánh giá và Thẩm định các Yêu cầu Phi chức năng}

Nhóm đã sử dụng công cụ kiểm thử tải Apache JMeter và Lighthouse để tiến hành đo lường các chỉ số phi chức năng thực tế của hệ thống trên môi trường sản phẩm:

\begin{table}[H]
\centering
\caption{Kết quả đo lường và thẩm định các yêu cầu phi chức năng (NFR Verification)}
\label{tab:nfr-verification}
\renewcommand{\arraystretch}{1.3}
\begin{tabular}{|p{3.2cm}|p{3.8cm}|p{4cm}|p{3.5cm}|}
\hline
\textbf{Chỉ số phi chức năng} & \textbf{Mục tiêu cam kết ban đầu} & \textbf{Kết quả đo lường thực tế} & \textbf{Đánh giá kết quả} \\
\hline
Thời gian phản hồi API thông thường & $< 300\text{ ms}$ & $68\text{ ms} - 124\text{ ms}$ & \textbf{Vượt cam kết} (Cực kỳ nhanh) \\
\hline
Thời gian tính toán gợi ý đối tác Jaccard & $< 1000\text{ ms}$ & $210\text{ ms} - 380\text{ ms}$ & \textbf{Vượt cam kết} (Nhờ tối ưu thuật toán tập hợp Set) \\
\hline
Thời gian phản hồi trợ lý ảo AI Gemini & $< 3.5\text{ s}$ & $1.8\text{ s} - 2.4\text{ s}$ & \textbf{Đạt cam kết} (Mô hình gemini-3.5-flash tối ưu) \\
\hline
Tỷ lệ chịu tải đồng thời & 50 requests/giây không có lỗi & 100 requests/giây, Error rate = 0.0\% & \textbf{Đạt cam kết} (Không xung đột dữ liệu) \\
\hline
Thời gian cập nhật trạng thái VietQR & $< 10\text{ s}$ sau chuyển khoản & $3.2\text{ s} - 4.5\text{ s}$ qua IPN Webhook & \textbf{Vượt cam kết} (Gần như tức thời) \\
\hline
Điểm đánh giá Lighthouse Performance & $\ge 85$ điểm & $94/100$ điểm trên thiết bị Desktop & \textbf{Vượt cam kết} (Tối ưu hóa chia chunk Vercel) \\
\hline
\end{tabular}
\end{table}

Các kết quả thực nghiệm định lượng ở trên minh chứng rằng hệ thống CoSpace không chỉ đáp ứng hoàn chỉnh toàn bộ các yêu cầu chức năng nghiệp vụ mà còn sở hữu chất lượng kỹ thuật xuất sắc về mặt hiệu năng, tính ổn định và bảo mật thông tin.